\begin{lemma}
\label{pre_pre_main}
For all \(n \in \N\), \(d \mid n\), prime powers \(q\), and \(f \in \CF_d (q)\),
\begin{equation}
\sum_{p \in \CF_n (q)}
\sum_{\substack{\alpha \in \F_{q^n}^\times \\ p(\alpha) = 0}}
\te(\alpha)^{\ell_f \cdot [n/d]_{q^d}}
=
\sum_{s \mid n} \moebius(n/s) (q^s-1) \delta_{q^s - 1 \mid \ell_f \cdot [n/d]_{q^d}}.
\end{equation}
\end{lemma}


\begin{proof}
By M\"obius inversion, it suffices to prove
\begin{equation}
\label{inverted}
\sum_{s \mid n}
\sum_{p \in \CF_s (q)}
\sum_{\substack{\alpha \in \F_{q^s}^\times \\ p(\alpha) = 0}}
\te (\alpha)^{\ell_f \cdot [n/d]_{q^d}}
=
(q^n - 1) \delta_{q^n - 1 \mid \ell_f \cdot [n/d]_{q^d}}.
\end{equation}
Corollary~\ref{theta_of_root_union} implies that, on the left side of~\eqref{inverted}, \(\te(\alpha)\) ranges precisely over all \((q^n-1)^\text{th}\) roots of unity.
The sum of the \((\ell_f \cdot [n/d]_{q^d})^\text{th}\) powers of all \((q^n-1)^\text{th}\) roots of unity is zero unless \(q^n-1\) divides \(\ell_f \cdot [n/d]_{q^d}\), in which case the sum is \(q^n-1\).
\end{proof}


\begin{lemma}
\label{post_pre_pre_main}
For all \(n \in \N, \, d \mid n\), prime powers \(q\), and \(b \mid q^n-1\),
\begin{equation}
\# \{f \in \CF_d (q) : \, \, b \mid \ell_f \cdot [n/d]_{q^d}\}
=
\frac{1}{d}
\sum_{c \mid d}
\moebius(d/c)
\frac{q^n-1}{\lcm \big (
[n/c]_{q^c}, b
\big )}
\end{equation}
\end{lemma}


\begin{proof}
By M\"obius inversion, it suffices to prove
\begin{equation}
\label{inverted_2}
\sum_{c \mid d} c \cdot \# \{f \in \CF_c (q) : \, \, b \mid \ell_f \cdot [n/c]_{q^c} \}
=
\frac{q^n-1}{\lcm \big (
[n/d]_{q^d}, b
\big )}.
\end{equation}
View each value \(\ell_f \cdot [n/c]_{q^c}\) as an element of \(\Z / (q^n-1)\), as in the context of Corollary~\ref{te_n_map_property}.
Modulo \(q^n-1\), there are exactly \(c\) distinct choices for each \(\ell_f \cdot [n/c]_{q^c}\).
Namely, given a choice for \(\ell_f\),
\[
\ell_f , q\ell_f , q^2\ell_f , \ldots, q^{c-1}\ell_f \quad \mod q^n-1
\]
are also valid choices for \(\ell_f\), and so
\[
\ell_f \cdot [n/c]_{q^c}, q\ell_f \cdot [n/c]_{q^c} , q^2 \ell_f \cdot [n/c]_{q^c} , \ldots, q^{c-1}\ell_f \cdot [n/c]_{q^c} \quad \mod q^n-1
\]
are all the possible choices for \(\ell_f \cdot [n/c]_{q^c}\) in \(\Z / (q^n-1)\).
Recalling definition~\eqref{te_n_def}, we see that those values are precisely the images under \(\te_n\) of the roots of \(f\).
Thus, another way to interpret the sum on the left side of \eqref{inverted_2} is
\[
\sum_{c \mid d} \sum_{f \in \CF_c (q)} \sum_{\substack{\alpha \in \F_{q^c} \\ f(\alpha) = 0}} \delta_{b \mid \te_n(\alpha)}.
\]
Therefore, by Corollary~\ref{te_n_map_property}, the left side of \eqref{inverted_2} counts the elements of \(\Z / (q^n-1)\) that are divisible by both \([n/d]_{q^d}\) and \(b\),
which is exactly the right side of \eqref{inverted_2}.
\end{proof}




\begin{proof}[Proof of Corollary~\ref{pre_main}]
Using Lemma~\ref{pre_pre_main} and expanding the \(k^\text{th}\) power in the statement, it remains to show
\[
\sum_{s_1,\ldots,s_k \mid n}
\# \{f \in \CF_d (q) : \lcm(q^{s_1}-1,\ldots,q^{s_k}-1) \mid \ell_f \cdot [n/d]_{q^d}\}
\cdot
\prod_{i=1}^k \moebius(n/s_i)(q^{s_i}-1)
\]
equals
\[
\frac{1}{d} \sum_{c \mid d} \moebius(d/c) C_{n,k,c} (q).
\]
By the definition of \(C_{n,k,c} (q)\),
this is equivalent to showing that
\[
\# \{f \in \CF_d (q) : \lcm(q^{s_1}-1,\ldots,q^{s_k}-1) \mid \ell_f \cdot [n/d]_{q^d} \}
\]
equals
\[
\frac{1}{d}
\sum_{c \mid d}
\moebius(d/c)
\frac{q^n-1}{\lcm \big ( [n/c]_{q^c}, q^{s_1}-1, \ldots, q^{s_k}-1 \big )}.
\]
This follows from Lemma~\ref{post_pre_pre_main}.
\end{proof}





































